\section{Complete finite certificates and their encoding}
\label{sec:global52}
The degree-three projector criterion admits a global certificate
hierarchy, with unknown sections included among its variables. This
uses classical real algebraic geometry, not a new Positivstellensatz.
We give the complete encoding because a variable count by itself is
not a complexity theorem.

\begin{proposition}[Description and local dual dimensions]
\label{prop:counts52}
Put $s=|X|$, $a=|\calA|$, and let $n$ be the variable count in
\eqref{eq:variables51}. The local criterion can be written with
\begin{equation}\label{eq:counts52}
 \begin{split}
 e={}&2\sum_{t=0}^N K_t^2+s(1+K_0+D)\\
     &+a\sum_{t=0}^{N-1}(K_t+K_tK_{t+1}+DK_t)
 \end{split}
\end{equation}
scalar equalities of degree at most three and $2n$ interval
inequalities. Redundant symmetry equations are counted, and there
is no projection-rank constraint. For a single transition of shape
$k\times\ell$, with projectors and mean matrices fixed, the complete
linear system $B\operatorname{vec}T=b$ has
\[
 B\in\R^{(k+k\ell+kD)\times k\ell},\qquad
 b,y\in\R^{k+k\ell+kD}.
\]
Its rows are, in order, row sums, projector-range invariance, and
observable intertwining. An infeasibility witness may be normalized
by $B^Ty\ge0$, $b^Ty=-1$.
\end{proposition}
\begin{proof}
Count the full symmetry and idempotence matrices, the seed mass,
range and mean equations, and the three transition equation groups.
For each variable use $[0,1]$ for probabilities or transition entries
and $[-1,1]$ for projector entries or means. This gives exactly $2n$
scalar bound inequalities. Symmetry and idempotence allow any rank;
unreachable protected directions are permitted by the equivalence.
For fixed geometric data the transition conditions are linear, with
$k$, $k\ell$, and $kD$ entries respectively. Equality-form Farkas
separation proves the dual alternative, and a negative $b^Ty$ can be
scaled to $-1$. Redundancies do not affect the alternative.
\end{proof}

For rational input, numerator and denominator bit lengths are included
in the description. For real-algebraic input, one may specify each
coefficient by an integer minimal polynomial and a rational isolating
interval, with primitive-element conversion if desired. The expanded
matrix formula has polynomial arithmetic description length in the
listed horizon, widths, dimension and alphabet; this does not bound
algebraic elimination time or primitive-element conversion by a
polynomial. A solution can be represented by real-algebraic
coordinates. Arbitrary real matrices satisfy the mathematical
criterion, but are not thereby finite effective inputs.

\begin{theorem}[A complete global residual hierarchy]\label{thm:sos52}
For rational or real-algebraic input, enumerate all scalar equality
residuals of Theorem~\ref{thm:local51} as $h_1,\ldots,h_e$, retaining
projector-range invariance. Let $Q$ be their variable box and define
\[
 q(x)=\sum_{j=1}^e h_j(x)^2,\qquad m_* =\min_{x\in Q}q(x).
\]
Then $\deg q\le6$, and the prescribed profile is feasible if and
only if $m_*=0$. Write the interval bounds as
$g_i(x)=(x_i-a_i)(b_i-x_i)\ge0$, and include
$g_0=n-\sum_i x_i^2\ge0$. The semidefinite bounds obtained from
\begin{equation}\label{eq:sos52}
 q-\lambda=\sigma_{-1}+\sum_{i=0}^n\sigma_i g_i,
 \qquad \sigma_i\text{ sums of squares},
\end{equation}
with increasing degree limits are nondecreasing and have supremum
$m_*$. Infeasibility is equivalent to the existence of one such
finite identity with $\lambda>0$. With algebraic input and rational
$0<\lambda<m_*$, the identity can be represented with real-algebraic
Gram matrices and verified exactly.
\end{theorem}
\begin{proof}
Compactness gives attainment of $m_*$. A sum of squares of residuals
vanishes precisely when every equality holds; the box contains all
remaining inequalities. This proves the feasibility assertion.
The added ball polynomial is redundant on $Q\subset[-1,1]^n$ and
makes the generated quadratic module Archimedean. Every identity
\eqref{eq:sos52} gives $q\ge\lambda$ on $Q$, hence $\lambda\le m_*$.
For every $\lambda<m_*$, the polynomial $q-\lambda$ is strictly
positive on $Q$. Putinar's theorem gives \eqref{eq:sos52} at a finite
degree \cite{Putinar93}; this is the standard convergent polynomial
optimization hierarchy \cite{Lasserre01}. It proves convergence and
finite detection of strict infeasibility, without a claim of finite
convergence at feasible boundary points.

At a fixed degree, coefficient matching and positive-semidefiniteness
of the Gram matrices are a finite semialgebraic system over the
real-algebraic coefficients. A real solution therefore has a
real-algebraic solution by real closed field transfer and sampling
\cite{BPR}. Coefficient identities and nonnegativity of all principal
minors give finite exact verification. In particular, a floating-point
semidefinite solver's status is not itself a certificate.
\end{proof}

The hierarchy optimizes over the full unknown projector and mean
variables, unlike the supplied-section linear alternative. It is a
global certificate procedure, not a polynomial-time width optimizer.
Degree growth, conditioning of numerical searches, and coefficient
heights remain substantive computational costs. The classical
hypotheses and convergence mechanism are explicitly credited; their
application does not convert this realization problem into a convex
problem in the original variables. The certificate checker accompanying
the article verifies supplied rational Gram identities and rejects
coefficient, positivity, normalization and range-invariance mutations.
It does not implement a general semidefinite search engine.
